% ===================================================================
%  नकसा नं. १ — इसकूल, माध्यमिक इसकूल, कक्षा।
%
%  Ceci n'est PAS une transcription : c'est une traduction, faite en
%  2026, en bhojpouri d'aujourd'hui. D'ou trois differences avec
%  texto/io et texto/fr, et trois seulement :
%
%   * pas de \nl ni de \cc — la traduction ne suit aucune ligne
%     imprimee, et les lignes de ce fichier ne sont que du confort ;
%   * pas de \begin{VUpage} — elle n'a ni feuillet ni folio ;
%   * le gras se pose sur le TERME seul, ponctuation dehors.
%
%  Les cles « %%K » sont celles de texto/io, macro pour macro pour
%  l'apparat, et les renvois \textsuperscript{(n)} visent les MEMES
%  objets de la planche, dans le MEME ordre.
%
%  LA SEULE COLONNE DU RELEVE DONT LA VOISINE NIE QU'ELLE EXISTE.
%  Quarante-six colonnes ont ete ecrites avant celle-ci et toutes se
%  defendaient contre une langue qui les tenait pour des langues : le
%  marathi contre le hindi, le gujarati contre le hindi, le persan
%  contre l'ourdou, le levantin contre l'arabe standard. Le bhojpouri
%  est administre en Inde comme un « dialecte du hindi », ne figure
%  dans aucune annexe officielle, et compte plus de cinquante
%  millions de locuteurs. La colonne n'est donc pas seulement voisine
%  de sa voisine : elle est censee lui appartenir.
%
%  IL N'Y A AUCUNE DEFENSE POSSIBLE AU CARACTERE, ET C'EST LA
%  PREMIERE FOIS. Le persan surveillait ی contre ي, le haoussa ses
%  quatre lettres crochues, le gujarati un signe devanagari glisse
%  dans son alphabet, le levantin quatre lettres persanes. Ici
%  l'ecriture est la meme devanagari, sans un signe de difference. Ce
%  qui separe les deux langues est ailleurs, et c'est le VERBE : la
%  copule बा / बानी / बाड़ें contre है / हैं, la negation
%  existentielle नइखे contre नहीं है, le passe en -ल — कइल, गइल,
%  आइल — contre le passe hindi en -आ. Les vingt-cinq regles de
%  kolonoj.py visent donc la conjugaison d'abord, les pronoms
%  ensuite, le lexique en dernier, et c'est l'inverse de l'ordre de
%  toutes les colonnes precedentes.
%
%  MAIS LE MOT EMPRUNTE, LUI, SE VOIT. Le bhojpouri ne supporte pas
%  qu'un mot commence par s + consonne : il pose une voyelle devant.
%  L'ecole n'est pas स्कूल mais « इसकूल », la gare n'est pas स्टेशन
%  mais इसटेसन. LE TITRE DE CE TABLEAU-CI EST DONC DEJA UNE PREUVE DE
%  LANGUE, et c'est le mot anglais qui la donne, pas le mot indien.
%  Second marqueur, dans le corps du tableau : le bhojpouri ecrit ब
%  la ou le hindi ecrit व — बियाकरन la grammaire (55), बिज्ञान la
%  science (11), बनस्पति la botanique (b).
%
%  parigi.py A ETE LU AVANT D'ECRIRE, ET NON APRES. Le bhojpouri est
%  a tete finale comme le gujarati : le determinant precede le
%  determine, le complement precede le verbe, et « sidas sur stulo en
%  sua katedro » se retourne mot pour mot. La colonne gujaratie avait
%  paye ce tableau-ci de deux fautes d'ordre pour avoir consulte
%  l'outil trop tard. Les six paires que parigi.py annonce pour le
%  tableau 1 ont ete posees sous les yeux avant la premiere ligne :
%
%      t01-01-1   (7) ← sur ← (8)      (8) ← en ← (6)
%      t01-02-1   (15) ← vitrizita ← (1)   (1) ← kun ← (2)
%      t01-05-2   (46) ← kun ← (45)
%      t01-08-1   (70) ← por ← (69)
%      t01-15-1   (61) ← di ← (59)
%      t01-16-1   (65) ← de ← (64)
%
%  LE VASISTAS (78) EST « रोशनदान », ET C'EST LA MEME REPONSE QUE LA
%  COLONNE GUJARATIE. Quarante-six colonnes, quarante-six reponses, et
%  celle-ci repete exactement la quarante-cinquieme : le donneur de
%  lumiere, mot persan (روزن‌دان) que le persan lui-meme n'emploie pas
%  — il dit دریچه, la petite porte. Le mot est parti pour l'Inde et il
%  y est reste ; le releve tient maintenant la preuve qu'il ne s'est
%  pas arrete au Gujarat mais a traverse jusqu'au Gange. UNE REPONSE
%  QUI SE REPETE N'EST PAS UNE REPONSE PERDUE : elle dit ou passe la
%  frontiere de la chose.
%
%  ET « मेज » LA TABLE (18) EST PORTUGAIS ICI AUSSI. La colonne
%  gujaratie l'avait releve la premiere — mesa, apporte par Diu et
%  Daman. Le mot est monte du littoral jusqu'a la plaine gangetique
%  sans que personne le tienne pour etranger, et l'ecolier de Chapra
%  qui pose sa कापी dessus ecrit du portugais sans le savoir. Au meme
%  alinea, कुरसी vient de l'arabe, बेंच de l'anglais et पाटी (42) est
%  du vieux fonds : quatre meubles d'une meme salle, quatre langues.
% ===================================================================

%%K t01-apar-0 apar
\VUsaut{2.0mm}
\VUcentre{12.6pt}{120}{समझावन-पोथी}
\VUsaut{3.2mm}
\VUcentre{8.4pt}{160}{खातिर}
\VUsaut{2.6mm}
\VUtitre{15.6pt}{0}{देल्मास सहायक नकसा}
\VUsaut{3.4mm}
\VUornamento{9pt}
\VUsaut{5.0mm}
\VUcentre{13.2pt}{90}{पहिल शृंखला}
\VUsaut{3.0mm}
\VUfilet{16mm}
\VUsaut{5.6mm}

%%K t01-apar-1 apar
\VUtitre{15.0pt}{60}{नकसा नं. १}
\VUsaut{1.4mm}
\VUtitre{11.4pt}{0}{इसकूल, माध्यमिक इसकूल, कक्षा।}
\VUsaut{2.2mm}
\VUfilet{20mm}
\VUsaut{6.4mm}

%%K t01-tit-1 sub
\VUpk{10.2pt}{40}{आम बयान।}
\VUsaut{3.0mm}

%%K t01-01-1 p
१. --- ई नकसा एगो \VUgras{माध्यमिक इसकूल} के \VUgras{कक्षा} देखावेला।
एह कक्षा में आठ गो \VUgras{मेज} \textsuperscript{(18)} आ आठ गो
\VUgras{बेंच} \textsuperscript{(32)} बा। हर बेंच पर तीन गो लइका बइठेलें।
\VUgras{मास्टर साहेब} \textsuperscript{(7)} आपन \VUgras{कुरसी}
\textsuperscript{(8)} पर, आपन \VUgras{मंच} \textsuperscript{(6)} पर बइठल
बाड़ें।

%%K t01-02-1 p
२. --- मंच के बाईं ओर काँच वाला \VUgras{आलमारी} \textsuperscript{(15)}
ठाढ़ बा। दाहिने ओर \VUgras{करिया पटिया} \textsuperscript{(1)} बा, आपन
\VUgras{झाड़न} \textsuperscript{(2)} आ \VUgras{खड़िया} \textsuperscript{(3)}
के साथे।

%%K t01-02-2 p
मंच के ऊपर \VUgras{भीत के नकसा} \textsuperscript{(9, 11, 12)} आ एगो
\VUgras{नकसा} \textsuperscript{(10)} लटकल बा।

%%K t01-03-1 p
३. --- सब लइका आपन \VUgras{जगहा} \textsuperscript{(17)} पर नइखन बइठल।
योहानेस \textsuperscript{(4)} करिया पटिया लगे ठाढ़ बा ; ऊ \VUgras{झाड़न}
हाथ में लेके \VUgras{रेखागणित के आकृति} मेटावत बा।

%%K t01-03-2 p
पेत्रुस मंच लगे बा ; ऊ \VUgras{लिखल काम} \textsuperscript{(5)} बटोरत बा
आ मास्टर साहेब लगे ले जात बा।

%%K t01-04-1 p
४. --- \VUgras{ई} मास्टर साहेब \VUgras{जिंदा भाषा} के मास्टर बाड़ें। ऊ
लइकन के परदेसी भाषा \VUgras{(जर्मन, अंगरेजी, स्पेनिस, इटालियन, रूसी भा
फ्रेंच)} बोले, पढ़े आ लिखे सिखावेलें।

%%K t01-05-1 p
५. --- कक्षा बहुते चाकर आ बहुते अँजोर बा। छोर पर दू गो
\VUgras{रोशनदान} \textsuperscript{(78)} वाला चाकर \VUgras{खिड़की} आ हवा
आ अँजोर खातिर \VUgras{काँच के केवाड़} बा।

%%K t01-05-2 p
ई कक्षा जाड़ नइखे। एकरा के गरम करे खातिर बीच में बहुते नीक
\VUgras{अँगीठी} \textsuperscript{(46)} बा, आपन \VUgras{नली}
\textsuperscript{(45)} के साथे।

%%K t01-05-3 p
भीत में \VUgras{लुगा के खूँटी} \textsuperscript{(58)} गाड़ल बा ; ओह पर
लइकन के टोपी आ कोट लटकल बा।

%%K t01-tit-2 sub
\VUsaut{5.2mm}
\VUpk{10.2pt}{40}{खेले के अँगना}
\VUsaut{3.6mm}

%%K t01-06-1 p
६. --- \VUgras{घड़ी} \textsuperscript{(63)} नौ बजके पाँच मिनट देखावत बा।

%%K t01-06-2 p
\VUgras{छुट्टी के बेरा} \textsuperscript{(76)} अबहीं खतम भइल आ पढ़ाई फेरु
शुरू होत बा। खेले के जगहा \VUgras{अँगना} \textsuperscript{(75)} में बा। हम
कुछ लइकन के खेलत देखत बानी। एगो \VUgras{निगरान मास्टर}
\textsuperscript{(74)} ओह लोग पर नजर राखत बाड़ें।

%%K t01-07-1 p
७. --- अँगना में हमनी के आउर कई गो लोग के देखत बानी जा, माने
\VUgras{इंस्पेक्टर} \textsuperscript{(73)}, \VUgras{हेडमास्टर}
\textsuperscript{(71)} आ \VUgras{उप-हेडमास्टर} \textsuperscript{(72)}।

%%K t01-07-2 p
इंस्पेक्टर इसकूलन के जाँच करेलें, हेडमास्टर माध्यमिक इसकूल चलावेलें, आ
उप-हेडमास्टर पढ़ाई पर नजर राखेलें।

%%K t01-07-3 p
चउथा मनई \VUgras{दरबान} \textsuperscript{(77)} बाड़ें। ऊ गैरहाजिर लोग के
नाँव लिखे खातिर बगल में हाजिरी-रजिस्टर दबवले बाड़ें।

%%K t01-08-1 p
८. --- अँगना के छोर पर \VUgras{कसरत के सामान} \textsuperscript{(70)} आ
\VUgras{हाथ के काम} \textsuperscript{(69)} खातिर कारखाना बा।

%%K t01-08-2 p
एह नकसा के दाहिने ओर हमनी के \VUgras{संगीत} आ \VUgras{चित्रकारी} के
\VUgras{कक्षा} देखे लागत बानी जा।

%%K t01-noto-1 noto
\VUnotes{143.2mm}{%
(*) \textit{बपतिसमा के नाँव} खातिर सलाह बा कि ओकनी के लैटिन रूप में
लिखल जाव। अनगिनत फरक से बाँचे के एहे एगो उपाय बा। ना त
\textit{Giovanni, Jean, Johann, Jan, Hans, John, Ivan} में से कवन
चुनल जाव ? का हमनी के बपतिसमा के नाँव खातिर अलगे सबदकोस बनाइब जा ?
लैटिन से ओकनी के कर्ता-रूप लीं आ कहीं : \VUgras{Ioannes, Ioanna;
Iulius, Iulia; Iakobus; Andreas; Lukas.} (Gram. Kompleta)।%
}

%%K t01-tit-3 sub
\VUpk{10.2pt}{40}{बिस्तार से बयान।}
\VUsaut{2.4mm}

%%K t01-tit-4 sub
\VUpk{10.2pt}{40}{नीक लइका।}
\VUsaut{3.0mm}

%%K t01-09-1 p
९. --- मंच के दाहिने ओर के पहिला बेंच के लइका \VUgras{हेनरिकुस} (*),
\VUgras{स्तेफानुस} आ \VUgras{कारोलुस} बाड़ें।

%%K t01-09-2 p
हेनरिकुस बहुते धेयान देवे वाला आ मेहनती बा ; ऊ मास्टर साहेब के सुनत बा।

ओकर मेज पर \VUgras{कापी} \textsuperscript{(28)}, \VUgras{किताब}
\textsuperscript{(29)}, \VUgras{शब्दकोश} \textsuperscript{(30)} आ
\VUgras{कलमदान} \textsuperscript{(31)} बा। ओकर किताब में बहुते
\VUgras{पन्ना} बा ; हर अध्याय में कई गो \VUgras{पैरा} आ हर
\VUgras{लाइन} में बहुते \VUgras{सबद} बा।

%%K t01-09-4 p
ओकर पड़ोसी स्तेफानुस \textsuperscript{(24)} लगनशील बा। ऊ आपन कापी में
अगिला बेर के पाठ लिखत बा। ओकर \VUgras{अच्छर} सुंदर बा। ओकर किताब बंद
बा ; हमनी के खाली ओकर \VUgras{जिल्द} \textsuperscript{(27)} देखत बानी
जा। ओकरा लगे ओकर \VUgras{परकार} \textsuperscript{(26)} धइल बा।

%%K t01-09-5 p
कारोलुस \textsuperscript{(23)} आपन किताब हाथ में पकड़ले बा ; ऊ आपन पाठ
फेरु पढ़े खातिर ओकरा के खोलत बा। हर लइका नियर ओकरो आगे \VUgras{दवात}
\textsuperscript{(25)} बा। ऊ कहल मानेला।

%%K t01-10-1 p
१०. --- बगल के मेज पर \VUgras{लुदोविकुस, अंतोनियुस} आ \VUgras{पाउलुस}
बाड़ें। लुदोविकुस आपन \VUgras{पाठ} \textsuperscript{(22)} सुनावत बा आ
मास्टर साहेब के \VUgras{सवालन} के जवाब देत बा। अंतोनियुस
\textsuperscript{(21)} बहुते बेधेयान बा, कबो-कबो त मास्टर साहेब ओकरा के
सजा देलें। पाउलुस \textsuperscript{(20)} नीक संगी बा, मिलनसार आ मददगार।
ऊ बहुते धेयान देवेला।

%%K t01-tit-5 sub
\VUsaut{4.2mm}
\VUpk{10.2pt}{40}{खराब लइका।}
\VUsaut{3.2mm}

%%K t01-11-1 p
११. --- हेनरिकुस के पाछे \VUgras{गेओर्गियुस} \textsuperscript{(38)} बा। ऊ
खराब लइका नइखे, बाकिर ऊ \VUgras{बतकुच्चन}, बेधेयान आ चंचल बा। ओकर
\VUgras{हलुकपन} आ \VUgras{बेअदबी} पढ़ाई में \VUgras{गड़बड़ी} पैदा करेला आ
अकसर ओकरा के कड़ा \VUgras{चेतावनी} मिलेला। ओकर \VUgras{कच्चा कापी}
\textsuperscript{(35)} मैल बा, ऊ आपन \VUgras{कलम} \textsuperscript{(34)}
से ओह पर \VUgras{सियाही के दाग} लगावेला, एहीसे ऊ हरमेसा आपन
\VUgras{रबर} \textsuperscript{(36)} राखेला ; ओकर \VUgras{बस्ता}
\textsuperscript{(33)} फाटल बा, काहेकि ऊ बहुते लापरवाह बा। ऊ आपन
\VUgras{पेंसिल के खोल} \textsuperscript{(37)} जिंदा भाषा के कक्षा में
काहे ले आवेला ?

%%K t01-11-2 p
ओकर पड़ोसी एदमुंदुस \textsuperscript{(39)} ओकरा के पसंद ना करे। ऊ सचहूँ
कुछ सुस्त बा, बाकिर चुप आ सांत बा। ऊ बतियावे ना ; बाकिर सुने के बदले ऊ
आपन \VUgras{पढ़े के किताब} के \VUgras{कहानी} पढ़ल पसंद करेला।

%%K t01-11-3 p
एह बेंच के तिसरका लइका \VUgras{लुदोविकुस} \textsuperscript{(40)} बा। ऊ
मेहनती बा। ऊ उजर \VUgras{कागज} पर लिखत बा। आपन आगे ऊ आपन
\VUgras{सोखता} \textsuperscript{(41)} धइले बा।

%%K t01-12-1 p
१२. --- मास्टर साहेब के बाईं ओर के दुसरका बेंच के लइका बहुते छोट बाड़ें।
पहिलका, छोटका \VUgras{एमिलियुस}, अबहीं ठीक से लिखे ना जानेला, ऊ आपन
\VUgras{पाटी} \textsuperscript{(42)}, आपन \VUgras{पाटी के पेंसिल} आ आपन
\VUgras{स्पंज} \textsuperscript{(43)} ले आवेला।

%%K t01-12-2 p
ओकरा बगल में \VUgras{आउगुस्तुस} आ \VUgras{बेर्नार्दुस}
\textsuperscript{(44)} बाड़ें। ई दुसरका नीक काम करेला, बाकिर ओकर
\VUgras{पहिनावा} बहुते बेढंग बा।

%%K t01-13-1 p
१३. --- ओह लोग के पाछे तीन गो \VUgras{अलसी} बाड़ें। \VUgras{आल्बेर्तुस}
आपन पाठ ना याद करे। \VUgras{याकोबुस} \textsuperscript{(47)} सुने के बदले
सूतेला। ओकर \VUgras{अलसीपन} से ओकरा \VUgras{डाँट} आ \VUgras{सजा} तक
मिलेला। आ \VUgras{क्रिस्तियानुस} \textsuperscript{(48)} बहुते गैरजिम्मेदार
बा। ऊ खराब \VUgras{बेवहार} करेला आ सुधरे के कवनो कोसिस नइखे करत।

%%K t01-14-1 p
१४. --- ओकर पड़ोसी लोग एकरा उलटा नीक चालचलन वाला बाड़ें।
\VUgras{एर्नेस्तुस} \textsuperscript{(51)} के ढंग आ नेम पसंद बा, ऊ हरमेसा
आपन \VUgras{पटरी} \textsuperscript{(50)} आ आपन \VUgras{पेंसिल}
\textsuperscript{(49)} इस्तेमाल करेला। ऊ \VUgras{बाहरी लइका} बा।

%%K t01-14-2 p
\VUgras{फ्रांसिस्कुस} \textsuperscript{(52)} \VUgras{होस्टल के लइका} बा,
ऊ वरदी पहिनेला, इसकूल में खाला आ इसकूल में सूतेला। ऊ नीक \VUgras{संगी}
बा।

%%K t01-14-3 p
माउरिकियुस आपन \VUgras{पेंसिल} आपन \VUgras{चाकू} \textsuperscript{(56)}
से छिलत बा, ऊ बहुते सलीकेदार बा।

%%K t01-14-4 p
ऊ आपन सब किताब ले आवेला, आपन \VUgras{बियाकरन} \textsuperscript{(55)},
आपन \VUgras{नोटबुक} \textsuperscript{(54)}, आ \VUgras{लिखे के गद्दी}
\textsuperscript{(53)} तक। ओकर \VUgras{इसकूल के झोरा} \textsuperscript{(57)}
ओकरा पाछे भुइयाँ पर बा।

%%K t01-14-5 p
कक्षा के छोर पर के दुनो मेज \VUgras{नीक लइकन} \textsuperscript{(19)} से
भरल बा।

%%K t01-tit-6 sub
\VUsaut{4.6mm}
\VUpk{10.2pt}{40}{चित्रकारी आ संगीत।}
\VUsaut{3.4mm}

%%K t01-15-1 p
१५. --- \VUgras{चित्रकारी के कक्षा} में भीत पर सुंदर \VUgras{नमूना}
टँगल बा आ छत से \VUgras{छाता} वाला \VUgras{बत्ती} \textsuperscript{(62)}
लटकल बा। \VUgras{मास्टर साहेब} \textsuperscript{(60)} बहुते धीरजवाला
बाड़ें, ऊ लइकन के \VUgras{चित्र} \textsuperscript{(61)} सुधारेलें आ ओह लोग
के देख के \VUgras{मूरत} \textsuperscript{(59)} बनावल सिखावेलें।

%%K t01-16-1 p
१६. --- \VUgras{संगीत के कक्षा} बहुते मजेदार लागत बा। ओहिजा संगीत
पढ़ल सिखावल जाला, \VUgras{सुर-रेखा} पर लिखल \VUgras{सुर}
\textsuperscript{(67)} (सा, रे, ग, म, प, ध, नी\,=\,C. D. E. F. G. A. B.)
गावल, \VUgras{कुंजी} पहिचानल आ मीठ \VUgras{धुन} गावल। संगीत के कागज
\VUgras{रिहल} \textsuperscript{(65)} पर धइल जाला। एगो लइका
\VUgras{पियानो बजावत} \textsuperscript{(64)} बा आ \VUgras{संगीत के
मास्टर} \textsuperscript{(66)} \VUgras{ताल देत} \textsuperscript{(68)}
बाड़ें।

%%K t01-17-1 p
१७. --- हमनी के \VUgras{हाथ के काम} \textsuperscript{(69)} वाला
\VUgras{कारखाना} के एगो कोना देखत बानी जा, जहाँ लइकन के लकड़ी छीलल आ
लोहा रेतल सिखावल जाला। ई सब इसकूलन में नइखे।

%%K t01-tit-7 sub
\VUsaut{5.0mm}
\VUpk{10.2pt}{40}{बिज्ञान के कक्षा।}
\VUsaut{3.6mm}

%%K t01-18-1 p
१८. --- गनित के मास्टर साहेब लइकन के \VUgras{रेखागणित, अंकगणित,
हिसाब} आ \VUgras{मीटरी नाप} सिखावेलें। पटिया पर हमनी के रेखागणित के
खास आकृति देखत बानी जा : \VUgras{गोल} \textsuperscript{(a)},
\VUgras{चौकोर} \textsuperscript{(b)}, \VUgras{तिकोन} \textsuperscript{(c)},
\VUgras{रेखा} \textsuperscript{(d)}।

%%K t01-18-2 p
हमनी के एगो \VUgras{सवाल} भी देखत बानी जा ; ई ना त \VUgras{जोड़} बा, ना
\VUgras{घटाव}, ना \VUgras{गुना} — ई \VUgras{भाग} \textsuperscript{(e)} बा।

%%K t01-19-1 p
१९. --- मीटरी नाप के पटिया पर हमनी के देखत बानी जा : \VUgras{मीटर}
\textsuperscript{(a)}, \VUgras{तराजू} \textsuperscript{(b)} आ ओकर
\VUgras{बटखरा}, \VUgras{लीटर} \textsuperscript{(e)}, \VUgras{डेकालीटर}
\textsuperscript{(f)}, \VUgras{डेसीलीटर} आ \VUgras{घन मीटर}
\textsuperscript{(d)}।

%%K t01-19-2 p
लइकन \VUgras{प्रकृति-बिज्ञान} \textsuperscript{(11)}, जीव-बिज्ञान
\textsuperscript{(a)}, बनस्पति-बिज्ञान \textsuperscript{(b)},
भू-बिज्ञान \textsuperscript{(c)} आ \VUgras{इतिहास} \textsuperscript{(12)}
भी पढ़ेलें।

%%K t01-19-3 p
आलमारी में \VUgras{भौतिकी} \textsuperscript{(15)} आ \VUgras{रसायन}
\textsuperscript{(16)} के औजार बा।

%%K t01-19-4 p
आलमारी के ऊपर बा : \VUgras{गोला} \textsuperscript{(14)}, \VUgras{शंकु}
आ \VUgras{भूगोला} \textsuperscript{(13)}।

%%K t01-19-5 p
पटियन के ऊपर दुनिया के पाँच हिस्सा देखावे वाला \VUgras{नकसा} लटकल बा :
\VUgras{यूरोप} \textsuperscript{(g)}, \VUgras{एशिया} \textsuperscript{(e)},
\VUgras{अफ्रीका} \textsuperscript{(f)}, \VUgras{अमेरिका}
\textsuperscript{(ab)} आ \VUgras{ओसीनिया} \textsuperscript{(h)}, आ दू गो
बड़का महासागर, \VUgras{प्रशांत महासागर} \textsuperscript{(c)} आ
\VUgras{अटलांटिक महासागर} \textsuperscript{(d)}।
